\documentclass[11pt,a4paper]{article}
\usepackage[margin=24mm]{geometry}
\usepackage[T1]{fontenc}
\usepackage{lmodern}
\usepackage{amsmath,booktabs,longtable,graphicx,xurl,pdflscape}
\usepackage[hidelinks]{hyperref}
\usepackage{microtype}
\setlength{\parindent}{0pt}
\setlength{\parskip}{5pt}
\setlength{\emergencystretch}{3em}
\title{Bicoherence Development Supplement\\\large Interpretable AI/Human Music Detection Research}
\author{Experimental methods and reproducibility record}
\date{7 September 2026}
\begin{document}
\maketitle
\begin{abstract}
This supplement documents a candidate cross-frequency audio phenomenon through
coefficient-domain counterexamples, continuous-waveform controls and a frozen
real instrument-note pilot. Independent replay verifies the numerical pipeline.
The observed responses concern known interventions, not AI/Human accuracy.
Baseline observability and a retained negative response limit interpretation.
No external admission gate or bicoherence classifier has passed or been fitted.
The evidence cutoff is 7 September 2026, 17:39 Asia/Shanghai. At that cutoff,
the all-combination classification experiment was separately under audit
and is not reported here.
\end{abstract}
\input{../BICOHERENCE_PHENOMENON_THESIS_SECTION_EN.tex}
\clearpage
\begin{landscape}
\begin{figure}[p]
\centering
\includegraphics[width=\linewidth,height=0.79\textheight,keepaspectratio]{../results/bc_nsynth_bifrequency_presentation_v1/bifrequency_maps.png}
\caption{Descriptive bifrequency maps after independent acceptance of the NSynth
pilot. Top: baseline instrument-balanced mean squared bicoherence and paired
closed-minus-independent differences. Bottom: eligible notes out of 54, using
paired eligibility for the differences. Circles mark the target fixed before
measurement. Different cells can represent different covered subsets. This is
not a time-frequency spectrogram, significance map or AI/Human classifier.}
\label{fig:bc-bifrequency-development}
\end{figure}
\end{landscape}
\clearpage
\section{Artifact identity and evidence boundaries}
All paths in the methods are relative to the experiment root
\path{artifacts/audio_phenomena_expansion_20260907/}. The human-readable reports
are \path{BC_CONTINUOUS_DEVELOPMENT_RESULTS_EN.md} and
\path{BC_NSYNTH_DEVELOPMENT_RESULTS_EN.md}. They retain complete code, raw-output
and audit references. The raw note experiment contains 1,299 products totaling
1,459,355,556 bytes excluding its commit marker, with identical local/NFS copies.
The visualization has three products totaling 711,513 bytes. This supplement is
a separate versioned PDF; it does not replace the earlier main thesis PDF.

\begin{description}
\item[NSynth frozen draft SHA-256]
\path{7c3c98389ee8c349bb69502b5336ff653ab0910cec58d9f88f181a43265adb9a}
\item[NSynth raw COMMIT SHA-256]
\path{c28e0136dfa76f7b0658f04ceeebcb3ae69bf32ddd19f02f149019d3f4e895c3}
\item[Independent note-replay receipt SHA-256]
\path{5ac2335ade358056f9ab123189aeee2122c5b37c58959b59747c4a69e5314c59}
\item[Bifrequency presentation COMMIT SHA-256]
\path{44bad92abe3e3b4686fc0cb43c01e9e30515b5c2cc19f58cc26f3c0f1e1fee3a}
\end{description}

The prospective GuitarSet design is
\path{BC_GUITARSET_DEVELOPMENT_DESIGN_EN.md}. At this cutoff, its official source
download, full archive/audio validation and scoped overlap audit are unfinished.
The planned player-and-score-disjoint development and reserved intersections
have not been frozen or BC-measured. Successful offline tests of its draft tool
do not establish dataset integrity or scientific measurement performance.

\section{Literature pointers}
The literature establishes motivation and mathematical interpretation, not
validated accuracy on this project's music cohort.
\begin{enumerate}
\item AlBadawy, E.~A., Lyu, S., and Farid, H. (2019).
\emph{Detecting AI-Synthesized Speech Using Bispectral Analysis}.
CVPR Workshops, pp.~104--109.
\url{https://openaccess.thecvf.com/content_CVPRW_2019/html/Media_Forensics/AlBadawy_Detecting_AI-Synthesized_Speech_Using_Bispectral_Analysis_CVPRW_2019_paper.html}
\item Fackrell, J.~W.~A., and McLaughlin, S. (1996).
\emph{Detecting nonlinearities in speech sounds using the bicoherence}.
Proceedings of the Institute of Acoustics, 18(9), pp.~123--130.
\url{https://www.ioa.org.uk/system/files/proceedings/jwa_fackrell_s_mclaughlin_detecting_nonlinearities_in_speech_sounds_using_the_bicoherence.pdf}
The source pointer is retained; no locally downloaded PDF is claimed for this item.
\end{enumerate}
\end{document}
